%==============================================================================
\section{Applications --- Ứng dụng của học máy}
%==============================================================================

\begin{ynghia}[title={Nguồn}]
\tsurl{https://www.tutorialspoint.com/machine_learning/machine_learning_applications.htm}.
\end{ynghia}

\begin{dongco}[title={Tổng quan}]
Học máy đã trở thành công nghệ phổ biến và ảnh hưởng đến nhiều khía cạnh cuộc sống: kinh doanh, y tế, giải trí, \ldots
Học máy giúp ra quyết định và tìm các phương án giải bài toán, từ đó nâng hiệu quả công việc ở nhiều lĩnh vực.
\end{dongco}

\begin{ynghia}[title={Một số ứng dụng thành công}]
Nguồn liệt kê: chatbot, dịch ngôn ngữ, nhận diện khuôn mặt, hệ gợi ý, xe tự hành, phát hiện đối tượng, phân tích ảnh y tế, \ldots
\end{ynghia}

\subsection{Danh mục các lĩnh vực ứng dụng (theo nguồn)}
\begin{itemize}
  \item Image and Speech Recognition
  \item Natural Language Processing
  \item Finance Sector
  \item E-commerce and Retail
  \item Automotive Sector
  \item Computer Vision
  \item Manufacturing and Industries
  \item Healthcare Sector
\end{itemize}

\subsection{Image and Speech Recognition}

\begin{ynghia}[title={Nhận dạng ảnh và giọng nói}]
Nhận dạng ảnh và giọng nói là hai mảng mà học máy đã cải thiện đáng kể.
Thuật toán ML được dùng trong các ứng dụng như:
\begin{itemize}
  \item nhận diện khuôn mặt,
  \item phát hiện đối tượng (object detection),
  \item nhận dạng giọng nói (speech recognition),
\end{itemize}
nhằm nhận diện và phân loại ảnh/giọng nói một cách chính xác.
\end{ynghia}

\subsection{Natural Language Processing (NLP)}

\begin{dinhnghia}[title={NLP là gì?}]
Xử lý ngôn ngữ tự nhiên (NLP) là lĩnh vực của khoa học máy tính nghiên cứu tương tác giữa máy và con người bằng ngôn ngữ tự nhiên.
NLP dùng thuật toán học máy để nhận diện từ loại, cảm xúc và nhiều khía cạnh khác của văn bản; nó phân tích, hiểu và sinh ngôn ngữ người.
\end{dinhnghia}

\begin{ynghia}[title={NLP xuất hiện ở đâu?}]
NLP hiện diện trong nhiều ứng dụng trên Internet như phần mềm dịch, công cụ tìm kiếm, chatbot, phần mềm sửa ngữ pháp, trợ lý giọng nói, \ldots
\end{ynghia}

\begin{phuongphap}[title={Một số ứng dụng NLP (theo nguồn)}]
\begin{itemize}
  \item Sentiment analysis
  \item Speech synthesis
  \item Speech recognition
  \item Text classification
  \item Chatbots
  \item Language translation
  \item Caption generation
  \item Document summarization
  \item Question answering
  \item Autocomplete in search engines
\end{itemize}
\end{phuongphap}

\subsection{Finance Sector}

\begin{dongco}[title={ML trong tài chính}]
Vai trò của học máy trong tài chính là duy trì giao dịch an toàn.
Trong giao dịch, dữ liệu được chuyển thành thông tin phục vụ ra quyết định.
\end{dongco}

\subsubsection{1) Fraud Detection}
\begin{ynghia}[title={Phát hiện gian lận}]
Học máy được dùng rộng rãi để phát hiện gian lận.
Phát hiện gian lận là quá trình dùng mô hình ML để theo dõi giao dịch và hiểu pattern trong dataset nhằm nhận diện hoạt động gian lận/đáng ngờ.
\end{ynghia}

\begin{ynghia}[title={Lợi ích}]
Thuật toán ML có thể phân tích lượng dữ liệu giao dịch rất lớn để phát hiện pattern và anomaly, giúp ngăn thất thoát tài chính và bảo vệ khách hàng.
\end{ynghia}

\subsubsection{2) Algorithmic Trading}
\begin{ynghia}[title={Giao dịch thuật toán}]
Thuật toán học máy được dùng để tìm các pattern phức tạp trong dataset lớn nhằm phát hiện tín hiệu giao dịch mà con người khó nhận ra.
\end{ynghia}

\begin{phuongphap}[title={Các ứng dụng tài chính khác (theo nguồn)}]
\begin{itemize}
  \item phân tích và dự báo thị trường chứng khoán,
  \item đánh giá và quản trị rủi ro tín dụng,
  \item phân tích an ninh và tối ưu danh mục,
  \item định giá và quản trị tài sản.
\end{itemize}
\end{phuongphap}

\subsection{E-commerce and Retail}

\begin{dongco}[title={ML trong thương mại điện tử/bán lẻ}]
Học máy giúp cải thiện kinh doanh qua hệ gợi ý và quảng cáo nhắm mục tiêu, tăng trải nghiệm người dùng.
Học máy cũng giúp marketing dễ hơn bằng cách tự động hoá tác vụ lặp.
\end{dongco}

\subsubsection{1) Recommendation Systems}
\begin{ynghia}[title={Hệ gợi ý}]
Hệ gợi ý cung cấp đề xuất cá nhân hoá dựa trên hành vi, sở thích và tương tác trước đó của người dùng.
Thuật toán ML phân tích dữ liệu người dùng và tạo đề xuất cho sản phẩm, dịch vụ và nội dung.
\end{ynghia}

\subsubsection{2) Demand Forecasting}
\begin{ynghia}[title={Dự báo nhu cầu}]
Công ty dùng ML để hiểu nhu cầu tương lai của sản phẩm/dịch vụ dựa trên xu hướng thị trường, hành vi khách hàng và dữ liệu bán hàng lịch sử.
\end{ynghia}

\subsubsection{3) Customer Segmentation}
\begin{ynghia}[title={Phân khúc khách hàng}]
ML có thể phân nhóm khách hàng thành các nhóm có đặc điểm tương tự.
Mục đích là hiểu hành vi và nhắm mục tiêu trải nghiệm cá nhân hoá.
\end{ynghia}

\subsection{Automotive Sector}

\begin{dongco}[title={Xe tự hành}]
Một đổi mới lớn là xe tự hành (driverless vehicles) có thể cảm nhận môi trường, tự lái vượt chướng ngại mà không cần hỗ trợ người.
Nguồn nói: xe tự hành dùng ML để phân tích liên tục xung quanh và dự đoán các khả năng xảy ra.
\end{dongco}

\subsection{Computer Vision}

\begin{dinhnghia}[title={Computer vision trong ML}]
Computer vision là ứng dụng của học máy dùng thuật toán và mạng nơ-ron để dạy máy rút trích thông tin có ý nghĩa từ ảnh/video số.
Nguồn nêu computer vision được dùng trong nhận diện khuôn mặt, chẩn đoán bệnh từ MRI, và xe tự hành.
\end{dinhnghia}

\begin{phuongphap}[title={Một số tác vụ (theo nguồn)}]
\begin{itemize}
  \item Object detection and recognition
  \item Image classification and recognition
  \item Facial recognition
  \item Autonomous vehicles
  \item Object segmentation
  \item Image reconstruction
\end{itemize}
\end{phuongphap}

\subsection{Manufacturing and Industries}

\begin{dongco}[title={Bảo trì dự đoán (predictive maintenance)}]
ML được dùng để theo dõi điều kiện vận hành của máy móc.
Predictive maintenance giúp phát hiện lỗi/khuyết tật của thiết bị đang hoạt động để tránh dừng máy bất ngờ.
Phát hiện bất thường cũng giúp lập kế hoạch bảo trì định kỳ.
\end{dongco}

\begin{ynghia}[title={Predictive maintenance là gì?}]
Predictive maintenance là quá trình dùng thuật toán ML để dự đoán khi nào cần bảo trì một máy (ví dụ thiết bị trong nhà máy).
Phân tích dữ liệu từ cảm biến và nguồn khác giúp phát hiện pattern báo hiệu khả năng hỏng, để bảo trì trước khi máy “gãy”.
\end{ynghia}

\subsection{Healthcare Sector}

\begin{dongco}[title={ML trong y tế}]
ML có nhiều ứng dụng trong y tế.
Ví dụ: phân tích ảnh y khoa để phát hiện bệnh (như ung thư), hoặc dự đoán kết quả điều trị dựa trên tiền sử bệnh và yếu tố khác.
\end{dongco}

\begin{phuongphap}[title={Một số ứng dụng (theo nguồn)}]
\begin{enumerate}
  \item \textbf{Medical Imaging and Diagnostics}: phân tích pattern trong ảnh để chỉ ra sự hiện diện của bệnh.
  \item \textbf{Drug Discovery}: phân tích dataset lớn để dự đoán hoạt tính sinh học và tìm thuốc tiềm năng bằng phân tích cấu trúc hoá học.
  \item \textbf{Disease Diagnosis}: nhận diện một số loại bệnh (ví dụ: ung thư vú, suy tim, Alzheimer, viêm phổi) bằng thuật toán ML.
\end{enumerate}
\end{phuongphap}

\begin{luuy}[title={Kết luận của nguồn}]
Đây chỉ là vài ví dụ. Khi ML tiếp tục phát triển, ta kỳ vọng ML được dùng ở nhiều lĩnh vực hơn, nâng hiệu quả, độ chính xác và sự tiện lợi.
\end{luuy}